% !TEX program = pdflatex
% !TEX root = paper2_GRSL_letter.tex
%% =====================================================================
%% IEEE GEOSCIENCE AND REMOTE SENSING LETTERS (GRSL) – 5-page letter
%% "Sensor-family composition of the in-situ archive changes the
%%  apparent bias and RMSD of satellite soil-moisture products."
%%
%% This version is reframed as a SATELLITE-VALIDATION letter (not an
%% in-situ sensor paper), per GRSL scope: in-situ material must serve a
%% remote-sensing result. Key design choices:
%%   * Fig. 2 is the PRODUCT-LEVEL consequence (bias/RMSD stratification
%%     of ERA5-Land, SMAP L3_E, SMAP-SMI). This is the central result.
%%   * The CRNS attribution and the full family-composition table are
%%     demoted to the Supplement (Fig. S1, Table S1); only the pooled,
%%     clustered CRNS contrast stays in the main text.
%%   * SMAP L3_E / SMAP-SMI processing and the mean-state scope of the
%%     metrics are stated explicitly.
%%   * Claim level softened to "sensor-family-structured residual";
%%     the central claim is "the ISMN benchmark is not sensor-family
%%     neutral," not a laboratory calibration proof.
%%   * Compile with pdflatex (self-contained thebibliography).
%%     Uses IEEEtran.cls (bundled on Overleaf / IEEE author kit).
%% =====================================================================
\documentclass[journal]{IEEEtran}
\usepackage{amsmath,amssymb}
\usepackage{graphicx}
\usepackage{url}
\usepackage[hidelinks]{hyperref}

\graphicspath{%
  {paper2_GRSL_code_package_20260708/figures/main/}%
  {paper2_GRSL_code_package_20260708/figures/supplement/}%
  {figures/main/}{figures/supplement/}%
  {../figures/main/}{../figures/supplement/}}

\begin{document}

\title{Sensor-Family Composition of the In-Situ Archive Changes the
Apparent Bias and RMSD of Satellite Soil-Moisture Products}

\author{Chang-Hwan~Park%
\thanks{Manuscript submitted July 8, 2026. \emph{(Corresponding author:
Chang-Hwan Park.)}}%
\thanks{This work was supported by Institute of Information \&
communications Technology Planning \& Evaluation (IITP) under the
Artificial Intelligence Convergence Innovation Human Resources
Development (\mbox{IITP-2026-RS-2023-00255968}) grant funded by the Korea
government(MSIT)}%
\thanks{C.-H. Park is with the Department of Civil Systems Engineering,
Ajou University, Suwon, Republic of Korea (e-mail: cpark@ajou.ac.kr).}%
\thanks{Processing scripts and derived station-reference and CRNS-ISMN
matchup summary tables are archived on Zenodo
(doi:10.5281/zenodo.21263063) and mirrored in the public source-code
repository (v20260708):
\url{https://github.com/ecomm77/soil_moisture_organic_bias}.}}

\markboth{IEEE Geoscience and Remote Sensing Letters}%
{Park: Sensor-Family Composition
Changes Satellite Soil-Moisture Validation Metrics}

\maketitle

\begin{abstract}
Global satellite and reanalysis soil-moisture products are commonly
validated against the International Soil Moisture Network (ISMN), whose
reference-error structure is often treated as independent of sensor
manufacturer. This letter quantifies how that composition assumption
affects reported product metrics. Across 1{,}625 ISMN station-sensors
(2020–2023), $78.4\,\%$ belong to one dielectric family (Stevens
HydraProbe), whose mean residual is drier than the references while
Decagon/METER TDR records are wetter—a $0.17$~m$^3$\,m$^{-3}$ family
swing that exceeds the ERA5-Land/SMAP~L3\_E/SMAP-SMI disagreement.
Propagated into product metrics, the mean-state bias correction inferred
from ISMN changes by a factor of five for ERA5-Land and changes sign for
SMAP~L3\_E depending on which sensor family is used, and the SMAP-SMI
mean-state RMSD increases by $41\,\%$ on the non-HydraProbe subsample. A
geography-controlled paired-sensor test and a dielectric-independent
cosmic-ray-neutron cross-check, including a USA-excluded sensitivity,
indicate that the residual structure is associated with sensor family
after geography and network checks. A 2021 overpass-matched SMAP
temporal check also shows family structure in ubRMSE and temporal
correlation metrics.
ISMN-based validation metrics should be reported with sensor-family
stratification or reweighting.
\end{abstract}

\begin{IEEEkeywords}
Soil moisture, validation, SMAP, ERA5-Land, ISMN, in-situ reference,
sensor bias, cosmic-ray neutron sensor.
\end{IEEEkeywords}

\IEEEpeerreviewmaketitle

\section{Introduction}
\IEEEPARstart{T}{he} quantitative skill of every global satellite
soil-moisture retrieval and reanalysis product is reported as a
validation statistic—a network-mean bias, RMSD, or correlation—
computed against in-situ observations. In practice this reference is
supplied almost exclusively by the International Soil Moisture Network
(ISMN)~\cite{Dorigo2021}, which aggregates more than twenty national
and regional networks into a common reference archive and is used
downstream to quote a single number per candidate product, including
for SMAP~\cite{ONeill2021,ChanSMAP2018} and ERA5-Land
\cite{MunozSabater2021}. Community validation protocols emphasize
site representativeness, error characterization, and transparent
reporting of bias, ubRMSE, and correlation metrics
\cite{Colliander2017,Gruber2020}.

Aggregating the archive this way implicitly assumes that its error
model is \emph{manufacturer-neutral}: that the reference is
instrument-blind, so a product can be scored without regard to the
particular sensors that happen to populate the archive. That
assumption need not hold. Dielectric soil-moisture sensors from
different vendors operate at different frequencies (e.g., $50$~MHz for
the Stevens HydraProbe~\cite{Seyfried2004}) and use different
calibration curves. If the archive is dominated by one family and that
family carries a systematic near-surface offset, the offset does not
average out—it is inherited by the validation benchmark of every
product scored against the archive.

This letter evaluates that assumption as a remote-sensing validation
question: \emph{how much do the reported bias and RMSD of ERA5-Land,
SMAP~L3\_E, and the SMAP-derived SMI product change when the ISMN
reference is stratified by sensor-manufacturer family?} The analysis
quantifies the product-level sensitivity directly, controls for
deployment geography with co-located sensor pairs and network-clustered
inference, and adds a dielectric-independent cosmic-ray-neutron (CRNS)
cross-check. The inference is not a laboratory calibration claim; it is
a validation-methodology result showing that the ISMN benchmark is not
sensor-family neutral for these mean-state product metrics.

\section{Data and Methods}

\subsection{ISMN extraction and family assignment}
All \texttt{*\_sm\_*.stm} files in the ISMN archive used here
(31{,}201 files, 2015–2023) were parsed; we retained records dated
entirely in 2020–2023 with nominal depth-from $\le 0.15$~m and
$>180$ valid hourly observations per year, giving 1{,}625 unique
(network, station, sensor, depth) records across 23 networks on all
continents. Each sensor string was matched to a manufacturer family
(HydraProbe family, Decagon/METER TDR, Campbell CS6xx, ThetaProbe,
TRIME, other); strings containing \texttt{hydraprobe} or
\texttt{stevens} form one HydraProbe family (the same Stevens Water
Monitoring 50~MHz probe under several model numbers). The full
composition is given in Supplement Table~S1.

\subsection{Reference products and processing}
For each product we sampled each station by nearest valid cell and
formed the residual $(\mathrm{ISMN}-\mathrm{reference})$.
\emph{ERA5-Land}~\cite{MunozSabater2021}: layer-1 (0–7~cm) volumetric
soil water, per-cell temporal mean over 2020–2023 at $0.1^\circ$.
\emph{SMAP~L3\_E}~\cite{ONeill2021}: SPL3SMP\_E v006; the analysis used
182 available days from 2023, combining ascending and descending passes
and retaining only recommended-quality retrievals (retrieval-quality
flag) into a $9$-km EASE-Grid mean field, with a $0.15^\circ$ cell
rejection radius. \emph{SMAP-SMI}~\cite{Zeng2025SMI}: $0.1^\circ$
product, native 2020–2023, on a stratified 192-day sample. A
temporal-window control (2020–2023 vs.\ 2023-only ERA5-Land) leaves
the family ordering and the HydraProbe-family dry residual unchanged
(family change $\le\!0.012$~m$^3$\,m$^{-3}$; Supplement).

\emph{Scope of the metrics.} These are \emph{mean-state} validation
metrics—the bias and RMSD of multi-year mean fields—not
overpass-matched ubRMSE or anomaly correlation~\cite{Entekhabi2010,DraperDAS2012}.
This is the quantity represented by a single network-mean number per
product. To connect the result to standard Cal/Val language, we also
report a 2021 overpass-matched SMAP~L3 SM\_P temporal check in the
Supplement.

\subsection{Attribution and confounding controls}
Because HydraProbe-heavy networks (SCAN, SNOTEL, USCRN) also occupy
particular climates and soils, sensor family and deployment geography
are correlated. We separate them three ways. (i)~\emph{Co-located
pairs}: stations hosting both a HydraProbe-family and a non-HydraProbe
sensor in the same ISMN grid cell ($n=21$), compared with a paired
$t$-test; both members lie in the retained near-surface layer
(depth-from $\le0.15$~m), and the corresponding depth constraint is
described in the Supplement. (ii)~\emph{Network-clustered
inference}: a cluster bootstrap (5{,}000 resamples) that resamples ISMN
networks as the sampling unit, plus leave-one-network deletion, used as
the primary uncertainty rather than na\"ive per-station bootstrap.
(iii)~\emph{Dielectric-independent cross-check}: because product
residuals alone cannot separate an in-situ sensor effect from
product-specific error, we add a CRNS reference whose neutron signal is
independent of soil permittivity~\cite{Zreda2012}. Four archives
(COSMOS-USA~\cite{ZhengHESS2024}, COSMOS-Europe
\cite{BogenaCOSMOSEurope2022}, COSMOS-UK~\cite{Cooper2021COSMOSUK},
CosmOz Australia~\cite{Hawdon2014CosmOz}) were matched to ISMN within
$0.1^\circ$–$0.5^\circ$. For COSMOS-USA, the CRNS$-$ISMN residual is
assembled as the Zheng et al. CRNS$-$ERA5-Land bias minus the
station-matched ISMN$-$ERA5-Land residual from our $0.1^\circ$ overlay;
the non-USA archives use direct CRNS$-$probe or local matchup contrasts.
The pooled HydraProbe-minus-non-HydraProbe contrast and its clustered
bootstrap are reported here, with archive-level detail and a
USA-excluded sensitivity in the Supplement.

\section{Results}

\subsection{The validation archive is manufacturer-asymmetric}
Of the 1{,}625 station-sensors, 1{,}274 ($78.4\,\%$) are
HydraProbe-family; all other families combined are $<\!22\,\%$. Within
the same retained archive, the family mean residual against ERA5-Land is
$-0.097$~m$^3$\,m$^{-3}$ for HydraProbe-family station-sensors but
$+0.068$~m$^3$\,m$^{-3}$ for Decagon/METER TDR (Fig.~\ref{fig:asymmetry}).
Thus the validation archive is not only manufacturer-asymmetric in its
counts; it also contains family-structured mean-state residuals. A
single archive-mean validation number is therefore influenced by which
families and sub-networks dominate the retained ISMN sample.

\begin{figure*}[t]
\centering
\includegraphics[width=0.92\textwidth]{paper_fig3_AB.png}
\caption{The ISMN validation archive is manufacturer-asymmetric.
(a)~Retained 2020–2023 near-surface ISMN station-sensors by
manufacturer family. HydraProbe-family records supply $78.4\,\%$ of the
archive. (b)~Mean ISMN$-$ERA5-Land residual by family; residual-panel
counts reflect valid ERA5-Land overlay coverage, and TRIME is omitted
because $n=2$. The HydraProbe-to-Decagon/METER contrast is
$0.17$~m$^3$\,m$^{-3}$, so archive composition can change a reported
product bias even before any retrieval algorithm changes.}
\label{fig:asymmetry}
\end{figure*}

\subsection{Product validation metrics depend on sensor-family
composition}
Stratifying the ISMN reference by family changes the validation
statistics of all three products (Fig.~\ref{fig:downstream}). The mean
residual swings from $-0.097$~m$^3$\,m$^{-3}$ for the HydraProbe family
to $+0.068$~m$^3$\,m$^{-3}$ for Decagon/METER TDR against ERA5-Land
(a $0.17$~m$^3$\,m$^{-3}$ swing, larger than the between-product
disagreement). Propagated into commonly reported validation metrics:
the ISMN-implied ERA5-Land mean-state bias correction is $-0.081$ over
the full archive, $-0.097$ on the HydraProbe subsample, and $-0.019$ on
the non-HydraProbe subsample—a factor-of-five change in amplitude.
For SMAP~L3\_E the same triple is $-0.016/-0.025/+0.008$: the inferred
correction \emph{flips sign}, so whether SMAP reads wetter or drier than
the in-situ reference depends on sensor-family composition. The mean-state
RMSD is likewise composition-dependent; for SMAP-SMI it rises from
$0.099$ (HydraProbe subsample) to $0.140$~m$^3$\,m$^{-3}$
(non-HydraProbe subsample), a $41\,\%$ inflation. Because HydraProbe
dominates the archive roughly three-to-one, the full-archive number is
weighted toward the HydraProbe-family subset.

\begin{figure*}[t]
\centering
\includegraphics[width=0.95\textwidth]{downstream_figure.png}
\caption{Sensor-family composition changes satellite-product
validation metrics. (a)~Mean-state ISMN$-$reference residual by family
for ERA5-Land, SMAP~L3\_E, and SMAP-SMI; the HydraProbe$-$Decagon/METER
swing ($\approx\!0.17$~m$^3$\,m$^{-3}$) exceeds the between-product
disagreement. (b)~Mean-state RMSD for each product under
the full archive and the HydraProbe-only and non-HydraProbe-only
subsamples; the SMAP-SMI RMSD inflates $41\,\%$ on the non-HydraProbe
subsample. (c)~Implied ISMN$-$reference bias correction by subsample;
the ERA5-Land amplitude changes by $\sim\!5\times$ and the SMAP~L3\_E
correction changes sign. (d)~Independent family-contrast checks:
co-located pairs give a non-HydraProbe-minus-HydraProbe difference of
$+0.075$~m$^3$\,m$^{-3}$ $[+0.031,+0.119]$, and the CRNS cross-check
gives a HydraProbe-minus-non-HydraProbe contrast of
$+0.155$~m$^3$\,m$^{-3}$ $[+0.042,+0.244]$ under network-cluster
bootstrap; site- and archive-cluster CIs also exclude zero. This is the
practical consequence for satellite soil-moisture validation and is
independent of the sensor-versus-deployment attribution.}
\label{fig:downstream}
\end{figure*}

Equal-weighting the HydraProbe and non-HydraProbe subsets changes the
ERA5-Land mean residual from $-0.081$ to $-0.058$~m$^3$\,m$^{-3}$ and
the SMAP~L3\_E residual from $-0.016$ to $-0.009$~m$^3$\,m$^{-3}$.
Equal weighting across the five non-TRIME families gives
$-0.038$~m$^3$\,m$^{-3}$ for ERA5-Land and $+0.001$~m$^3$\,m$^{-3}$ for
SMAP~L3\_E, illustrating how a simple reweighting protocol converts the
diagnostic into an actionable reporting rule. The two primary
independent checks remain significant after Holm correction
($p_{\mathrm{Holm}}=4.6\times10^{-4}$ for CRNS and
$4.7\times10^{-3}$ for co-located pairs).

\subsection{Overpass-matched temporal check}
A real 2021 SMAP~L3 SM\_P (AM) overpass-matched check adds the Cal/Val
metrics that the mean-state analysis intentionally does not target
(Supplement Table~S3, Fig.~S2). Across 785 station-sensors and 78{,}600
matched station-days, the pooled HydraProbe versus non-HydraProbe bias
contrast is small but significant ($-0.032$ versus
$-0.036$~m$^3$\,m$^{-3}$; $p=9.8\times10^{-3}$) because the
non-HydraProbe bin mixes wetter Decagon/METER TDR records
($+0.058$~m$^3$\,m$^{-3}$) with drier Campbell and other records
($-0.107$ and $-0.071$~m$^3$\,m$^{-3}$). The headline family pair
therefore remains aligned with the mean-state result: HydraProbe is
drier than Decagon/METER by $0.090$~m$^3$\,m$^{-3}$, and Decagon/METER
also has higher Pearson $R$ ($0.622$ versus $0.317$) and anomaly $R$
($0.423$ versus $0.329$). Pooled ubRMSE differs ($0.062$ versus
$0.057$~m$^3$\,m$^{-3}$; $p=8.2\times10^{-4}$), while pooled Pearson
$R$ does not ($p=0.70$) and pooled anomaly $R$ is marginal ($p=0.055$).
This single-year diagnostic does not replace a full multi-product
temporal validation, but it shows that family structure is visible in
standard satellite-validation metrics as well as in mean-state
bias/RMSD.

\subsection{The family residual is robust to geography and network checks}
Two controls address whether the effect reflects deployment geography.
At the $n=21$ co-located pairs—which hold location, climate, season,
and soil constant—the non-HydraProbe sensor reads wetter by
$+0.075$~m$^3$\,m$^{-3}$ (paired $t=+3.19$, $p=4.7\times10^{-3}$; CI
$[+0.031,+0.119]$). Under network-clustered resampling the
HydraProbe-family residual remains negative and its CI excludes zero,
and every leave-one-network deletion preserves the effect
($p<10^{-3}$). The supported interpretation is therefore limited but
specific: the residuals remain sensor-family structured after
network-level robustness checks.

\subsection{Dielectric-independent CRNS cross-check}
Pooling the 92 CRNS$\leftrightarrow$ISMN pairs that permit a
within-archive family contrast (58 HydraProbe, 34 non-HydraProbe), the
dielectric-free neutron reference sits closer to the non-HydraProbe
families: mean HydraProbe residual $+0.111$ versus non-HydraProbe
$-0.044$~m$^3$\,m$^{-3}$ (Welch's $t=+4.07$, $p=2.3\times10^{-4}$). The
contrast survives every leave-one-archive deletion ($p<10^{-3}$), every
leave-one-network deletion ($p<10^{-3}$), and a cluster bootstrap whose
CIs exclude zero under network, CRNS-site, and archive clustering. The
COSMOS-USA component is the synthesized CRNS$-$ERA5 minus ISMN$-$ERA5
residual described above; when COSMOS-USA is removed, the remaining
archives still preserve the contrast ($t=+3.91$,
$p=6.3\times10^{-4}$). The Australian archive ($n=5$) is directionally
consistent but too small for an independent test. Per-archive results
are given in Supplement Fig.~S1; the cross-check adds a
permittivity-independent line of evidence for the same family ordering
seen in Fig.~\ref{fig:downstream}.

\section{Discussion and Conclusion}
The global ISMN benchmark is dominated by a single dielectric family,
and the in-situ$-$reference residual is sensor-family structured. The
consequence for remote sensing is direct and does not
require resolving whether the offset is a calibration effect or a
deployment effect: because the HydraProbe-family dry residual is not
distributed symmetrically across the archive, the mean-state bias and
RMSD reported for ERA5-Land, SMAP~L3\_E, and SMAP-SMI depend
on which sensor family populates the reference—to the point of
reversing the sign of an inferred SMAP correction. Manufacturer
neutrality of the ISMN benchmark is therefore not a safe assumption for
these mean-state metrics.

These results support treating sensor family as an explicit
stratification, fixed-effect, or reweighting term in ISMN-based
validation, with global product skill reported by family or on a
family-balanced subsample; even equal HydraProbe/non-HydraProbe
weighting materially changes the reported ERA5-Land and SMAP~L3\_E
mean residuals. Two limitations bound the claim. First, the primary
metrics are mean-state summaries, although a real 2021 SMAP temporal
check shows family structure in ubRMSE and anomaly $R$ as well. Second,
the attribution, although supported by a geography-controlled pair test
and a dielectric-independent CRNS cross-check, remains a
residual-structure result rather than a laboratory calibration proof. A
sensitivity analysis controlling for SoilGrids~2.0 clay fraction and
bulk density~\cite{Poggio2021} yields the same family ordering,
indicating that the effect is not explained by soil texture alone.

\section*{Acknowledgment}
The author thanks the ISMN, COSMOS-Europe, COSMOS-UK, CosmOz, and NASA
NSIDC data providers.

\begin{thebibliography}{16}

\bibitem{Dorigo2021}
W.~Dorigo \emph{et al.}, "The International Soil Moisture Network:
serving Earth system science for over a decade," \emph{Hydrol. Earth
Syst. Sci.}, vol.~25, pp.~5749–5804, 2021,
doi:10.5194/hess-25-5749-2021.

\bibitem{Gruber2020}
A.~Gruber \emph{et al.}, "Validation practices for satellite soil
moisture retrievals: What are (the) errors?" \emph{Remote Sens.
Environ.}, vol.~244, Art.~111806, 2020,
doi:10.1016/j.rse.2020.111806.

\bibitem{MunozSabater2021}
J.~Mu\~{n}oz-Sabater \emph{et al.}, "ERA5-Land: a state-of-the-art
global reanalysis dataset for land applications," \emph{Earth Syst.
Sci. Data}, vol.~13, pp.~4349–4383, 2021,
doi:10.5194/essd-13-4349-2021.

\bibitem{ONeill2021}
P.~E.~O'Neill \emph{et al.}, "SMAP Enhanced L3 Radiometer Global and
Polar Grid Daily 9~km EASE-Grid Soil Moisture, Version~6," NASA NSIDC
DAAC, 2021, doi:10.5067/M20OXIZHY3RJ.

\bibitem{ChanSMAP2018}
S.~K.~Chan \emph{et al.}, "Development and assessment of the SMAP
enhanced passive soil moisture product," \emph{Remote Sens.
Environ.}, vol.~204, pp.~931–941, 2018,
doi:10.1016/j.rse.2017.08.025.

\bibitem{Colliander2017}
A.~Colliander \emph{et al.}, "Validation of SMAP surface soil moisture
products with core validation sites," \emph{Remote Sens. Environ.},
vol.~191, pp.~215–231, 2017, doi:10.1016/j.rse.2017.01.021.

\bibitem{Entekhabi2010}
D.~Entekhabi \emph{et al.}, "The Soil Moisture Active Passive (SMAP)
mission," \emph{Proc. IEEE}, vol.~98, pp.~704–716, 2010,
doi:10.1109/JPROC.2010.2043918.

\bibitem{DraperDAS2012}
C.~S.~Draper, R.~H.~Reichle, G.~J.~M. De Lannoy, and Q.~Liu,
"Assimilation of passive and active microwave soil moisture
retrievals," \emph{Geophys. Res. Lett.}, vol.~39, Art.~L04401, 2012,
doi:10.1029/2011GL050655.

\bibitem{Zeng2025SMI}
J.~Zeng, P.~Wang, H.~Ma, P.~Shi, and C.~Zhang, "Global daily soil
moisture products at $0.1^\circ$ and $0.25^\circ$," Aerosp. Inf. Res.
Inst., CAS, Zenodo record 15015557, 2025,
doi:10.5281/zenodo.15015557.

\bibitem{Seyfried2004}
M.~S.~Seyfried and M.~D.~Murdock, "Measurement of soil water content
with a 50-MHz soil dielectric sensor," \emph{Soil Sci. Soc. Am. J.},
vol.~68, pp.~394–403, 2004, doi:10.2136/sssaj2004.3940.

\bibitem{Zreda2012}
M.~Zreda \emph{et al.}, "COSMOS: the COsmic-ray Soil Moisture
Observing System," \emph{Hydrol. Earth Syst. Sci.}, vol.~16,
pp.~4079–4099, 2012, doi:10.5194/hess-16-4079-2012.

\bibitem{BogenaCOSMOSEurope2022}
H.~R.~Bogena \emph{et al.}, "COSMOS-Europe: a European network of
cosmic-ray neutron soil moisture sensors," \emph{Earth Syst. Sci.
Data}, vol.~14, pp.~1125–1151, 2022,
doi:10.5194/essd-14-1125-2022.

\bibitem{Cooper2021COSMOSUK}
H.~M.~Cooper \emph{et al.}, "Soil moisture estimates from the UK
COSMOS-UK network," \emph{Earth Syst. Sci. Data}, vol.~13,
pp.~1737–1757, 2021, doi:10.5194/essd-13-1737-2021.

\bibitem{Hawdon2014CosmOz}
A.~Hawdon, D.~McJannet, and J.~Wallace, "Calibration and correction
procedures for cosmic-ray neutron soil moisture probes located across
Australia," \emph{Water Resour. Res.}, vol.~50, pp.~5029–5043, 2014,
doi:10.1002/2013WR015138.

\bibitem{ZhengHESS2024}
Y.~Zheng \emph{et al.}, "Evaluation of reanalysis soil moisture
products using cosmic ray neutron sensor observations across the
globe," \emph{Hydrol. Earth Syst. Sci.}, vol.~28, pp.~1999–2022,
2024, doi:10.5194/hess-28-1999-2024.

\bibitem{Poggio2021}
L.~Poggio \emph{et al.}, "SoilGrids 2.0: producing soil information
for the globe with quantified spatial uncertainty," \emph{SOIL},
vol.~7, pp.~217–240, 2021, doi:10.5194/soil-7-217-2021.

\end{thebibliography}

\end{document}
